\section{珂朵莉树}

\subsection{珂朵莉树（区间赋值 + 区间第 k 小 + 区间幂和）}

\par \noindent 珂朵莉树的核心是用 \texttt{set} 维护若干值相同的连续段（区间）。\texttt{split(pos)} 将包含 \texttt{pos} 的区间分裂成 \texttt{[l, pos-1]} 和 \texttt{[pos, r]} 两段并返回后者。区间操作时先 split 左右端点，再遍历中间所有段进行处理。只有在存在大量\textbf{区间赋值}操作时才有优秀的均摊复杂度。

\begin{minted}{c++}
const int mod=1e9+7;

inline ll read()
{
    ll x=0,f=1;char c=getchar();
    while(c<'0'||c>'9'){if(c=='-')f=-1;c=getchar();}
    while(c>='0'&&c<='9'){x=(x<<3)+(x<<1)+(c^48);c=getchar();}
    return x*f;
}

ll n,m,seed,vmax; 
ll rnd()
{
    ll ret=seed;
    seed=(seed*7+13)%mod;
    return ret;
}
const int N=1e5+10;
ll a[N];
struct qwq
{
    int l,r;mutable ll v;
    bool operator<(const qwq&x)const{return l<x.l;}
};
set<qwq>s;
auto split(int pos)
{
    auto it=s.lower_bound((qwq){pos,0,0});
    if(it!=s.end()&&it->l==pos)return it;
    --it;
    int l=it->l,r=it->r;ll v=it->v;
    s.erase(it);s.insert((qwq){l,pos-1,v});
    return s.insert((qwq){pos,r,v}).first;
}

void add(int l,int r,ll x)
{
    auto itr=split(r+1),itl=split(l);
    for(auto it=itl;it!=itr;it++)it->v+=x;
}

void assign(int l,int r,ll x)
{
    auto itr=split(r+1),itl=split(l);
    s.erase(itl,itr);s.insert((qwq){l,r,x});
}

ll rk(int l,int r,int k)
{
    vector<pair<ll,int>>v;
    auto itr=split(r+1),itl=split(l);
    for(auto it=itl;it!=itr;it++)v.push_back({it->v,it->r-it->l+1});
    sort(v.begin(),v.end());
    for(auto x:v)
    {
        k-=x.second;
        if(k<=0)return x.first;
    }
}
ll ksm(ll a,ll b,ll p)
{
    a%=p;
    ll res=1;
    while(b)
    {
        if(b&1)res=res*a%p;
        b>>=1;
        a=a*a%p;
    }
    return res;
}

ll query(int l,int r,ll x,ll y)
{
    ll res=0;
    auto itr=split(r+1),itl=split(l);
    for(auto it=itl;it!=itr;it++)
        res=(res+(it->r-it->l+1)*ksm(it->v,x,y)%y)%y;
    return res;
}

signed main()
{
    n=read(),m=read(),seed=read(),vmax=read();
    for(int i=1;i<=n;i++)a[i]=(rnd()%vmax)+1,s.insert((qwq){i,i,a[i]});
    for(int i=1;i<=m;i++)
    {
        int op=(rnd()%4)+1,l=(rnd()%n)+1,r=(rnd()%n)+1;
        if(l>r)swap(l,r);
        if(op==1)
        {
            ll x=(rnd()%vmax)+1;
            add(l,r,x); 
        }
        if(op==2)
        {
            ll x=(rnd()%vmax)+1;
            assign(l,r,x);
        }
        if(op==3)
        {
            int x=(rnd()%(r-l+1))+1;
            cout<<rk(l,r,x)<<endl;
        }
        if(op==4)
        {
            ll x=(rnd()%vmax)+1,y=(rnd()%vmax)+1;
            cout<<query(l,r,x,y)<<endl;
        }
    }
    return 0;
}
\end{minted}

\par \noindent 说明：
\begin{itemize}
\item \textbf{split(pos)}：将包含 \texttt{pos} 的区间分裂为 \texttt{[l, pos-1]} 和 \texttt{[pos, r]}，返回后者的迭代器。
\item \textbf{assign(l, r, x)}：区间推平，是珂朵莉树的核心操作，也是保证复杂度的关键。先 split 出 \texttt{[l, r+1)}，再删除中间所有段，最后插入一整段。
\item \textbf{add / rk / query}：区间加、区间第 $k$ 小、区间幂和，都是先 split 出区间，再遍历所有段分别处理。
\item \textbf{mutable}：让 \texttt{set} 中的元素可以被直接修改（如 \texttt{it->v += x}），不需要重新插入。
\end{itemize}

\subsection{放弃珂朵莉树（CF1638E）}

\par \noindent 放弃珂朵莉树： (CF1638E，学习自jiangly) 维护同色段，可以直接用线段树，也可以用 std::map<int, Info>，其中 key 为段的左端点，左闭右开这样 next 就是右端点，value 为段的信息，这种方式非常好写。每次进行区间染色 $[l,r)$ 的时候，先把 $l,r$ 两点加入map，然后就可以直接循环删除中间的，然后处理修改。
\begin{tcolorbox}
\par \noindent 给定一个长度为 $n$ 的序列，初始时所有元素的值为 0 ，颜色为 1。你需要实现以下三种操作：
\begin{itemize}
\item Color $l$ $r$ $c$ ：把 $[l,r]$这段的元素颜色改为 $c$
\item Add $c$ $x$：把所有颜色为 $c$ 的元素值都加 $x$
\item Query $i$：输出元素 $i$ 的值
\end{itemize}
\end{tcolorbox}
\begin{minted}{c++}
using namespace std;
using ll=long long;

const int N=1000010;
int n,m;
// map<key,value>  key  表示区间左端点
//                 next 表示区间右端点 [it, next(it))
map<int,int>s;

map<int,int>::iterator split(int x)
{
    auto it=prev(s.upper_bound(x));

    if(it->first==x)
        return it;

    return s.emplace(x,it->second).first;
}
// 区间修改单点查询树状数组
namespace Fenwick
{
    ll t[N];
    void add(int k,ll x){for(;k<=n;k+=(k&-k)) t[k]+=x;}
    ll sum(int k){ll v=0;for(;k;k-=(k&-k)) v+=t[k];return v;}
    void update(int l,int r,ll x){add(l,x);add(r,-x);}
}
using namespace Fenwick;

ll tag[N]; // tag[i]颜色是i的懒标记

signed main()
{
    ios::sync_with_stdio(false);
    cin.tie(nullptr);
    cout.tie(nullptr);

    cin>>n>>m;
    // 初始只有一个区间 [1, n+1) 颜色是 1 维护的每一段区间颜色都是相同的
    s[1]=1;
    s[n+1]=0;

    char op[10];
    while(m--)
    {
        cin>>op;

        if(*op=='C')
        {
            int l,r,c;
            cin>>l>>r>>c;
            r++;// 找到 [l, r + 1) 区间
            map<int,int>::iterator it=split(l);
            map<int,int>::iterator ti=split(r);
            
            // 边打标机边合并(erase)区间 最终得到 [l, r + 1)
            for(;it!=ti;it=s.erase(it)) 
                update(it->first,next(it)->first,tag[it->second]-tag[c]);
            
            s[l]=c; 
        }
        else if(*op=='A')
        {
            int c,x;
            cin>>c>>x;
            tag[c]+=x;
        }
        else
        {
            int i;
            cin>>i;
            map<int,int>::iterator it=split(i);
            cout<<sum(i)+tag[it->second]<<'\n';
        }
    }

    return 0;
}
\end{minted}